\documentclass[a4paper,11pt]{ltjsarticle}
\usepackage{luatexja}
\usepackage{amsmath}
\usepackage{amssymb}

\begin{document}

\begin{figure}[h]
  \centering
  % 図6.5の画像はここに挿入
  \caption{水素分子 $\mathrm{H_2}$ のエネルギー準位}
\end{figure}

\subsection{6.2.2 ダイヤモンドにおける共有結合}

結晶での共有結合の典型例はダイヤモンドにおける結合である。ダイヤモンドは炭素（C）からなる結晶である。炭素原子の電子配置は
$(1s)^2(2s)^2(2p)^2$
である。このうち $1s$ 電子は内殻電子であるため，結合への寄与はほとんどなく，結合に寄与するのは $(2s)^2(2p)^2$ の価電子である。すでに2.3節で説明したように，ダイヤモンドの結晶構造は図6.6で示す構造をしており，配位数は4，つまり1つの炭素原子に注目したとき，その原子に隣接する原子の数は4つである。1つの炭素原子を中心とする4つの結合はどれも同等であることと，結合に寄与する価電子が $(2s)^2(2p)^2$ のように2種類存在するということとは合わない。これに対して，ポーリング\footnote{ライナス・ポーリング（Linus Pauling, 1901～1994）は共鳴の概念や電気陰性度の概念などで有名なアメリカの化学者である。1954年に「化学結合の本性と複雑な分子の構造研究」によりノーベル化学賞を受賞。ジェームズ・ワトソンの著書『二重らせん（The Double Helix）』にも登場するように，DNAの構造決定の競争におけるワトソン，クリックらの強力なライバルであった。また，1962年に地上核実験に対する反対運動の業績によりノーベル平和賞も受賞しており，人類史上2度ノーベル賞を受賞した数少ない人の中の1人である。}
は炭素原子が4つの等価な結合を形成することを説明するために，\textbf{混成軌道}（hybridized orbital）という考え方を導入した。その考え方とは，$2s$ 電子の1つを $2p$ 軌道に移して $(2s)^1(2p)^3$ という電子配置にし，さらにこの4つの軌道を混成させて4つの等価な軌道をつくるというものである。具体的には

\begin{align}
\phi_1 &= \frac{1}{2}(\chi_{2s} - \chi_{2p_x} - \chi_{2p_y} - \chi_{2p_z}) \tag{6.13}\\
\phi_2 &= \frac{1}{2}(\chi_{2s} - \chi_{2p_x} + \chi_{2p_y} + \chi_{2p_z}) \tag{6.14}\\
\phi_3 &= \frac{1}{2}(\chi_{2s} + \chi_{2p_x} - \chi_{2p_y} + \chi_{2p_z}) \tag{6.15}\\
\phi_4 &= \frac{1}{2}(\chi_{2s} + \chi_{2p_x} + \chi_{2p_y} - \chi_{2p_z}) \tag{6.16}
\end{align}

のように4つの軌道をつくる\footnote{$p_x, p_y, p_z$ は方位量子数1である波動関数の線形結合でつくられる電子状態であり，それぞれ$x$軸，$y$軸，$z$軸に対称な関数である。例えば式(6.13)で与えられる関数は $[x,y,z]=[-1,-1,-1]$ 方向を向いた形の関数となる。}。これらの軌道は1個の$s$軌道と3個の$p$軌道からつくられているので\textbf{$sp^3$混成軌道}（$sp^3$ hybridized orbital）という。式からわかるようにすべての軌道に1個の $2s$ 軌道と3個の $2p$ 軌道が同じ割合で含まれている。$p$軌道の波動関数の対称性を考慮すれば式(6.13)～式(6.16)はそれぞれ $[-1\ -1\ -1]$，$[-1\ 1\ 1]$，$[1\ -1\ 1]$，$[1\ 1\ -1]$ 方向を向いた形の関数を表す。つまり，図6.6の下側に示すような正四面体の頂点となる。これに対して，隣接した炭素原子については式(6.13)～式(6.16)中の$p$軌道に関する関数の正負を入れ替えることで，$[1\ 1\ 1]$，$[1\ -1\ -1]$，$[-1\ 1\ -1]$，$[-1\ -1\ 1]$ を向いた形の関数をつくることができる。したがって，隣接する原子どうしが互いに向かい合うような方向で電子を共有することで結合が形成される。ダイヤモンドにおける共有結合では，隣接する原子間でスピンが反対向きとなる2個の電子を共有する。

共有結合の特徴は，ダイヤモンドの場合に見られるように，結合をつくる原子間に電子が多く分布することである。またダイヤモンド中の2つの炭素原子間の結合エネルギーが7.3\,eVであることからもわかるように，結合エネルギーが大きいことも共有結合の特徴である。共有結合の他の例には Si や Ge などがある。Si 原子間の結合エネルギーは4.6\,eV，Geでは3.9\,eVであり，やはり大きな値である。

\begin{figure}[h]
  \centering
  % 図6.6の画像はここに挿入
  \caption{ダイヤモンドの結晶構造とダイヤモンドの結合における混成軌道の方向}
\end{figure}

\subsection{6.3 イオン結合}

結晶でのイオン結合（ionic bond）の典型例は塩化ナトリウム（NaCl）である。ナトリウム（Na）原子の電子配置は
\[
(1s)^2(2s)^2(2p)^6(3s)^1
\]
塩素（Cl）原子の電子配置は
\[
(1s)^2(2s)^2(2p)^6(3s)^2(3p)^5
\]
であり，それぞれ次のようにイオン化すると閉殻構造となる。
\[
\mathrm{Na^+}: (1s)^2(2s)^2(2p)^6
\]
\[
\mathrm{Cl^-}: (1s)^2(2s)^2(2p)^6(3s)^2(3p)^6
\]

このようにしてできたイオン間のクーロン引力によってイオン結合が生じる。イオン結合によって形成される結晶を\textbf{イオン結晶}（ionic crystal）という。

\subsubsection{6.3.1 マーデルング定数}

イオン結合の場合，イオンを点電荷として近似することで\footnote{点電荷でなくてもイオンが球対称であれば以下の考え方が成り立つ。}比較的容易にクーロン引力によって生じるポテンシャルエネルギー（これを\textbf{マーデルングエネルギー}（Madelung energy）という）を見積もることができる。図6.7に示すように，結晶中の$i$番目のイオンと$j$番目のイオンが距離 $r_{ij} = |\bm{r}_j - \bm{r}_i|$ だけ離れているとき，この2つのイオン間に働くクーロン力によって生じるポテンシャルエネルギーは
\begin{equation}
V_{ij} = \frac{Z_i Z_j e^2}{4\pi\varepsilon_0 r_{ij}} \tag{6.17}
\end{equation}
で与えられる。ここで，$Z_i e$ は $i$ 番目のイオンの電荷，$\varepsilon_0$ は真空の誘電率である。$i$ 番目のイオンと，それ以外のすべてのイオンとの間に働くクーロン力によって生じるポテンシャルエネルギーは，すべての組み合わせについて式(6.17)の総和を求めることによって
\begin{equation}
V_i = \sum_{j\neq i} V_{ij} = \frac{e^2}{4\pi\varepsilon_0}\sum_{j\neq i}\frac{Z_i Z_j}{r_{ij}} \tag{6.18}
\end{equation}
と得られる。

最近接イオンどうしのクーロン力によって生じるポテンシャルエネルギー $-\dfrac{Z_+Z_- e^2}{4\pi\varepsilon_0 d}$ で式(6.18)を割った
\begin{equation}
M = \frac{d}{Z_+Z_-}\sum_{j\neq i}\frac{Z_i Z_j}{r_{ij}} \tag{6.19}
\end{equation}
は\textbf{マーデルング定数}（Madelung constant）\footnote{4.9.3項でも紹介したマーデルング（Erwin Madelung）による。}と呼ばれ，結晶構造によってその値が定まる。ここで，$Z_+e$ と $Z_-e$ はそれぞれイオン結晶中の陽イオンと陰イオンの電荷であり，$d$は最近接イオン間距離である。

例えば，塩化ナトリウム構造のマーデルング定数を，図6.8に示すような，慣用単位胞（辺の長さが格子定数$a$に等しい立方体）を$3\times3\times3$倍した立方体構造に対して近似的に求めてみよう。この立方体の中心に位置する $\mathrm{Cl^-}$ イオンを座標原点に設定すると，この $\mathrm{Cl^-}$ イオンにとって最近接イオンは $(\pm\frac{a}{2},0,0)$，$(0,\pm\frac{a}{2},0)$，$(0,0,\pm\frac{a}{2})$ に位置する6個の $\mathrm{Na^+}$ イオンであり，いずれも原点から $r_1 = \frac{a}{2} = d$ だけ離れている。$Z_{\mathrm{Na}}=+1$，$Z_{\mathrm{Cl}}=-1$ であるから，最近接イオンによるマーデルング定数への寄与は
\[
6 \times \frac{d}{r_1} = 6
\]
である。

第二近接イオンは $(\pm\frac{a}{2},\pm\frac{a}{2},0)$，$(\pm\frac{a}{2},\mp\frac{a}{2},0)$，$(\pm\frac{a}{2},0,\pm\frac{a}{2})$，$(\pm\frac{a}{2},0,\mp\frac{a}{2})$，$(0,\pm\frac{a}{2},\pm\frac{a}{2})$，$(0,\pm\frac{a}{2},\mp\frac{a}{2})$ に位置する12個の $\mathrm{Cl^-}$ イオンであり，いずれも原点から $r_2 = \frac{\sqrt{2}a}{2} = \sqrt{2}d$ だけ離れている。$\mathrm{Cl^-}$ イオンどうしの組み合わせであることに注意すれば，マーデルング定数への寄与は
\[
-12 \times \frac{d}{r_2} = -\frac{12}{\sqrt{2}}
\]
である。

第三近接イオンは $(\pm\frac{a}{2},\pm\frac{a}{2},\pm\frac{a}{2})$，$(\pm\frac{a}{2},\pm\frac{a}{2},\mp\frac{a}{2})$，$(\pm\frac{a}{2},\mp\frac{a}{2},\pm\frac{a}{2})$，$(\mp\frac{a}{2},\pm\frac{a}{2},\pm\frac{a}{2})$ に位置する8個の $\mathrm{Na^+}$ イオンであり，いずれも原点から $r_3 = \frac{\sqrt{3}a}{2} = \sqrt{3}d$ だけ離れているので，マーデルング定数への寄与は
\[
8 \times \frac{d}{r_3} = \frac{8}{\sqrt{3}}
\]
である。

同様にすべての組み合わせについて総和を求める。ただし，立方体の面上に位置するイオンについてはその半分だけが立方体内に含まれるのでマーデルング定数への寄与を$\frac{1}{2}$とし，同様の考えから，立方体の辺上に位置するイオンについてはその寄与を$\frac{1}{4}$，立方体の頂点に位置するイオンについてはその寄与を$\frac{1}{8}$とすると，
\begin{align*}
M &\simeq 6 - \frac{12}{\sqrt{2}} + \frac{8}{\sqrt{3}} - \frac{6}{\sqrt{4}} + \frac{24}{\sqrt{5}} - \frac{24}{\sqrt{6}} - \frac{12}{\sqrt{8}} + \frac{24}{\sqrt{9}} - \frac{8}{\sqrt{12}}\\
&\quad + \frac{1}{2}\times\left(\frac{6}{\sqrt{9}} - \frac{24}{\sqrt{10}} + \frac{24}{\sqrt{11}} + \frac{24}{\sqrt{13}} - \frac{48}{\sqrt{14}} + \frac{24}{\sqrt{17}}\right)\\
&\quad + \frac{1}{4}\times\left(-\frac{12}{\sqrt{18}} + \frac{24}{\sqrt{19}} - \frac{24}{\sqrt{22}}\right) + \frac{1}{8}\times\frac{8}{\sqrt{27}}\\
&\simeq 1.747
\end{align*}
と近似値が求められる。実際には塩化ナトリウム構造のマーデルング定数は1.748と求められている\footnote{マーデルング定数は無理数なので，小数点以下4桁の近似値を求める際に，右のような球の半径をだんだん大きくしてイオンからの寄与を求める方法に基づく式 $\dfrac{8}{\sqrt{3}} - \dfrac{6}{\sqrt{4}} + \dfrac{24}{\sqrt{5}} \cdots$ が教科書に載っていることがあるが，“On the convergence of lattice sums and Madelung's constant”, D. Borwein, J. M. Borwein and K. F. Taylor, J. Math. Phys., 26, 2999 (1985) によって，この方法では正しい値が求められないことが示されており，正しく求められるマーデルング定数を計算するには，陽イオンと陰イオンからの寄与を打ち消すようにするなどの工夫が重要となる。}。

参考として，塩化ナトリウム構造を含むさまざまな結晶構造のマーデルング定数を表6.1に示す。マーデルング定数の値が1より大きいということは，結晶中において，同種イオンの間に働く斥力によってポテンシャルエネルギーが増加するにもかかわらず，1個の陽イオンと1個の陰イオンが距離$d$だけ離れているときのポテンシャルエネルギー $-\dfrac{e^2}{4\pi\varepsilon_0 d}$ よりもエネルギーが低下することを意味する。つまり，陽イオンと陰イオンにとっては，分子よりも結晶となった方が結合エネルギーが大きくさらに安定となる。

\subsubsection{6.3.2 塩化ナトリウムの結合エネルギー}

塩化ナトリウムの格子定数が $a = 0.564\,\mathrm{nm}$，マーデルング定数が $M = 1.748$ であることを用いるとマーデルングエネルギーは
\begin{equation}
V_i = -\frac{Me^2}{2\pi\varepsilon_0 a} = 1.43\times10^{-18}\,\mathrm{J} = 8.92\,\mathrm{eV}
\end{equation}
と見積もられる。

実際の結晶においては，イオンどうしが近づくと斥力も働くようになるため，図6.9に示すように，ポテンシャルエネルギーはマーデルング定数を用いた見積もりの90\%程度の値になる。実験から得られているエネルギーは8.14\,eVであるから，マーデルング定数を用いた見積もりがかなり良いことがわかる。

ナトリウム（Na）原子と塩素（Cl）原子から塩化ナトリウム（NaCl）結晶が生成する反応についてのエネルギー収支は以下のようになる。

\begin{align*}
\mathrm{Na} + 5.1\,\mathrm{eV}\ \text{（イオン化エネルギー）} &\rightarrow \mathrm{Na^+} + (-e)\ \text{（電子）}\\
(-e)\ \text{（電子）} + \mathrm{Cl} &\rightarrow \mathrm{Cl^-} + 3.6\,\mathrm{eV}\ \text{（電子親和力）}\\
\mathrm{Na^+} + \mathrm{Cl^-} &\rightarrow \mathrm{NaCl}\ \text{（結晶）} + 8.1\,\mathrm{eV}\ \text{（格子エネルギー）}
\end{align*}

ここで，\textbf{イオン化エネルギー}（ionization energy）は原子から電子を引き離してイオンにするのに必要なエネルギーであり，\textbf{電子親和力}（electron affinity）は原子に電子を付け加えたときに放出あるいは吸収されるエネルギーである。エネルギーが放出される場合，電子親和力は負であると定義する。したがって，NaとClのそれぞれがイオンとなるときには，Naのイオン化エネルギーからClの電子親和力を差し引いた
\begin{equation}
5.1\,\mathrm{eV} - 3.6\,\mathrm{eV} = 1.5\,\mathrm{eV}
\end{equation}
だけエネルギー的に損をするが，総合的にはバラバラの原子でいるよりもイオン結合で結晶になった方が
\begin{equation}
8.1\,\mathrm{eV} + 3.6\,\mathrm{eV} - 5.1\,\mathrm{eV} = 6.6\,\mathrm{eV}
\end{equation}
だけエネルギー的に安定となる。つまり，塩化ナトリウムの結合エネルギーは6.6\,eVである。イオン結合の一般的な特徴としては，結合エネルギーが大きいこと，電子分布が原子に集中していることなどがあげられる。

\begin{figure}[h]
  \centering
  % 図6.8の画像はここに挿入
  \caption{塩化ナトリウム構造のマーデルング定数を近似的に求めるために考える立方体構造。ここでは図を見やすくするためにイオンの大きさを小さくしてある。}
\end{figure}

\begin{table}[h]
  \centering
  \caption{さまざまな結晶構造のマーデルング定数}
  \begin{tabular}{lc}
    \hline
    結晶構造 & マーデルング定数 \\
    \hline
    閃亜鉛鉱構造 & 1.638 \\
    ウルツ鉱構造 & 1.641 \\
    塩化ナトリウム構造 & 1.748 \\
    塩化セシウム構造 & 1.763 \\
    蛍石型構造 & 2.519 \\
    \hline
  \end{tabular}
\end{table}

\begin{figure}[h]
  \centering
  % 図6.9の画像はここに挿入
  \caption{イオン結合におけるポテンシャルエネルギーのイオン間距離依存性}
\end{figure}

\end{document}